\documentclass[12pt,a4paper]{article}
\usepackage[utf8]{inputenc}
\usepackage[T1]{fontenc}
\usepackage{amsmath}
\usepackage[left=2.00cm, right=2.00cm, top=2.00cm, bottom=2.00cm]{geometry}
\usepackage{amssymb}
\usepackage[italian]{babel}

\title{Soluzioni Tutorato 2\\Complementi di Matematica per le Scienze Chimiche\\ A.A. 2021/2022}
\author{prof. Emanuele Bottazzi}
\date{}

\usepackage{amsthm}
\theoremstyle{definition}
\newtheorem{exe}{Esercizio}
\newtheorem*{sol}{Soluzione}

\begin{document}
\maketitle

\begin{exe}
    Insieme di definizione...
\end{exe}
\begin{sol}
    a) $|y| \leq |x|$; b) $x^2 + y^2 < 1$; c) $|x| < 1$ e $x^2 + y^2 \leq 4$.
\end{sol}

\begin{exe}
    Derivate parziali...
\end{exe}
\begin{sol}
    a) $\frac{\partial f}{\partial x} = \frac{\frac{y}{x^2+y^2} \cdot (x^2+y^2) - \arctan(x/y) \cdot 2x}{(x^2+y^2)^2}$;
    b) $\frac{\partial f}{\partial x} = \frac{-y \sin(xy)}{1 + \cos(xy)}$.
\end{sol}

\begin{exe}
    Ortogonalità gradiente...
\end{exe}
\begin{sol}
    $\nabla f(0,1) = (1, -4)$. $\mathbf{u} = (\sqrt{2}, \frac{1}{2\sqrt{2}})$.
    $\nabla f(0,1) \cdot \mathbf{u} = \sqrt{2} - \frac{4}{2\sqrt{2}} = \sqrt{2} - \sqrt{2} = 0$.
\end{sol}

\begin{exe}
    Derivata direzionale...
\end{exe}
\begin{sol}
    a) $\nabla f(1/2, -\sqrt{3}/2) = (1, -\sqrt{3})$. $v = \frac{1}{\sqrt{3}}(\sqrt{2}, 1)$. $D_v f = \frac{\sqrt{2} - \sqrt{3}}{\sqrt{3}}$.
\end{sol}

\begin{exe}
    Piano tangente...
\end{exe}
\begin{sol}
    a) $z = 3 + 2(x-2) + (y-1) = 2x + y - 2$.
\end{sol}

\begin{center}
    prof. Emanuele Bottazzi\\
    emanuele.bottazzi@unipv.it
\end{center}

\end{document}
